\documentclass[12pt]{article}
\usepackage{amsmath,amssymb,amsfonts}
\usepackage[margin=1in]{geometry}

\begin{document}

\section*{Chapter 3, Problem 18}

\textbf{Problem.} Show that there exists no nonsingular matrix $P$ of order 2 such that $P^{-1}AP = D$, where $D$ is a diagonal matrix and
\[
A = \begin{pmatrix} 1 & 0 \\ 1 & 1 \end{pmatrix}.
\]
This example shows that not every matrix can be diagonalized by a similarity transformation.

\bigskip
\textbf{Solution.}

\textbf{Eigenvalues of $A$:} The characteristic equation is
\[
\det(A - \lambda I) = (1 - \lambda)^2 = 0.
\]
So $\lambda = 1$ is the only eigenvalue, with algebraic multiplicity 2.

\medskip
\textbf{Eigenvectors:} Solve $(A - I)x = 0$:
\[
A - I = \begin{pmatrix} 0 & 0 \\ 1 & 0 \end{pmatrix}, \qquad \begin{pmatrix} 0 & 0 \\ 1 & 0 \end{pmatrix}\begin{pmatrix} x_1 \\ x_2 \end{pmatrix} = 0 \quad \Rightarrow \quad x_1 = 0.
\]
The eigenspace is spanned by $\begin{pmatrix} 0 \\ 1 \end{pmatrix}$ alone. The geometric multiplicity is 1.

\medskip
\textbf{Why diagonalization fails:} If $A$ were diagonalizable via $P^{-1}AP = D$, then $D$ would have the eigenvalues of $A$ on its diagonal, so $D = I$. But then $A = PIP^{-1} = I$, which contradicts the fact that $A \neq I$ (since $A_{21} = 1$).

Alternatively, diagonalization requires $n$ linearly independent eigenvectors to form the columns of $P$. Here $n = 2$ but the eigenspace is only 1-dimensional, so we cannot find 2 linearly independent eigenvectors. Therefore no diagonalizing $P$ exists. \hfill $\blacksquare$

\end{document}
